\section{The Last Supper: Anticipation Without Fiction}

Calling the Supper an anticipation can sound weak if it means merely a dramatic
forecast. Catholic doctrine means more. Christ did not offer the apostles a
symbolic rehearsal whose truth would begin only the next afternoon. He gave
them his Body and Blood and instituted the sacramental form in which his one
sacrifice would remain with the Church.

The temporal asymmetry must nevertheless be retained. At the Supper the
historical shedding of blood was still future. On the Cross it occurred. The
Council of Trent and later magisterial teaching present the Supper as the
institution and true sacramental offering of Christ's Body and Blood before
their bloody immolation on the Cross; the Second Vatican Council says that
Christ instituted the Eucharistic sacrifice on the night he was betrayed in
order to perpetuate the sacrifice of the Cross (Trent, session XXII,
chs.~1--2; \emph{Sacrosanctum
Concilium} 47; CCC 1323, 1365--1367).

``Anticipatory sacramental presence'' is therefore a careful theological
description if its limits are stated. It says that the coming self-offering
gave the Supper its sacrificial content and that Christ truly gave himself
under sacramental signs. It does not say that the bloody sacrifice had already
happened historically, that the apostles watched the Crucifixion from the
upper room, or that Christ died twice. Supper and Cross are historically
distinct moments within Christ's one self-offering: at the first he institutes
the Eucharist and truly offers and gives his Body and Blood under sacramental
signs; at the second he consummates that self-offering through bloody
immolation.

